% 图 51.7 —— 由 `checks/emit_hand.py` 自动拼出（probe 精测 + prims2 图元分解）
%%
% ⚠ 这是**稿子不是成品**：跑 `checks/checkhand.py 51.7` 看指标 + 看叠合图。
%%
%% ⚠ 待核：
%   · 本图 probe 报出阴影族 1 个 —— **本脚本不生成阴影**（要照原书手写）；prims2 的 strokes 里可能含阴影线，核对时留意别画重
%   · 可疑字形 2 项（空心圈 ⊙ 常被认成字母 o）—— 见文件末
%%
\documentclass[12pt]{standalone}
\usepackage{fontspec}
\setsansfont{Arial}
\newfontfamily\mfallback{DejaVu Sans}
\usepackage{tikz}
\begin{document}
\begin{tikzpicture}[x=1pt, y=-1pt, line width=3.0pt, line cap=round, line join=round]

  %% ════ 填充区（prims2 的 fills）════
  \fill (1085.0,182.0) -- (1091.0,183.0) -- (1094.0,187.0) -- (1091.0,193.0) -- (1085.0,194.0) -- (1089.0,195.0) -- (1096.0,193.0) -- (1098.0,191.0) -- (1099.0,186.0) -- (1095.0,182.0) -- cycle;
  \fill (260.0,465.0) -- (254.0,467.0) -- (254.0,473.0) -- (257.0,474.0) -- (259.0,477.0) -- (265.0,473.0) -- (263.0,470.0) -- (264.0,468.0) -- cycle;

  %% ════ 直线段（prims2 的 strokes）════
  \draw[line width=3.0pt] (200.0,469.0) -- (185.0,469.0);
  \draw[line width=6.0pt] (109.0,416.0) -- (219.0,301.0);
  \draw[line width=9.4pt] (219.0,299.0) -- (211.0,297.0);
  \draw[line width=5.0pt] (343.0,421.0) -- (342.0,410.0);
  \draw[line width=7.0pt] (264.0,326.0) -- (272.0,321.0);
  \draw[line width=3.0pt] (172.0,238.0) -- (186.0,238.0);
  \draw[line width=4.9pt] (107.0,419.0) -- (110.8,421.8);
  \draw[line width=7.0pt] (110.8,421.8) -- (177.0,422.0);
  \draw[line width=4.0pt] (344.0,422.0) -- (421.0,421.0);
  \draw[line width=9.5pt] (335.0,186.0) -- (345.0,173.0);
  \draw[line width=6.1pt] (447.0,420.0) -- (442.0,423.0);
  \draw[line width=4.0pt] (226.0,421.0) -- (342.0,422.0);
  \draw[line width=5.0pt] (586.0,509.0) -- (595.0,508.0);
  \draw[line width=5.7pt] (264.0,327.0) -- (266.0,332.0);
  \draw[line width=2.0pt] (201.0,476.0) -- (185.0,476.0);
  \draw[line width=6.0pt] (272.0,321.0) -- (332.9,197.1);
  \draw[line width=7.0pt] (332.9,197.1) -- (335.0,187.0);
  \draw[line width=7.0pt] (272.0,321.0) -- (272.0,331.0);
  \draw[line width=4.0pt] (564.0,508.0) -- (554.0,508.0);
  \draw[line width=6.0pt] (975.0,227.0) -- (982.0,225.0);
  \draw[line width=9.7pt] (174.0,413.0) -- (179.0,421.0);
  \draw[line width=7.0pt] (756.0,136.0) -- (744.0,129.0);
  \draw[line width=5.7pt] (230.0,240.0) -- (235.0,238.0);
  \draw[line width=5.0pt] (528.0,507.0) -- (528.0,502.0);
  \draw[line width=4.0pt] (267.0,332.0) -- (271.0,332.0);
  \draw[line width=4.0pt] (405.1,52.0) -- (415.7,51.0);
  \draw[line width=4.0pt] (695.8,109.0) -- (742.0,128.0);
  \draw[line width=6.0pt] (266.0,333.0) -- (226.0,420.0);
  \draw[line width=4.0pt] (103.0,173.0) -- (84.0,173.0);
  \draw[line width=5.0pt] (742.0,119.0) -- (744.0,128.0);
  \draw[line width=4.0pt] (104.0,174.0) -- (104.0,412.2);
  \draw[line width=5.0pt] (104.0,412.2) -- (107.0,416.0);
  \draw[line width=4.0pt] (104.0,172.0) -- (105.0,142.0);
  \draw[line width=6.0pt] (334.0,186.0) -- (325.7,190.3);
  \draw[line width=6.0pt] (325.7,190.3) -- (222.0,299.0);
  \draw[line width=4.0pt] (105.0,173.0) -- (133.0,174.0);
  \draw[line width=8.7pt] (984.0,223.0) -- (984.0,215.0);
  \draw[line width=5.0pt] (104.0,421.0) -- (18.0,423.0);
  \draw[line width=5.0pt] (648.0,508.0) -- (649.0,502.0);
  \draw[line width=4.0pt] (105.0,422.0) -- (104.0,473.0);
  \draw[line width=4.5pt] (343.0,437.0) -- (343.0,423.0);
  \draw[line width=7.0pt] (225.0,421.0) -- (179.0,421.0);
  \draw[line width=8.8pt] (173.0,430.0) -- (178.0,423.0);
  \draw[line width=5.5pt] (263.0,327.0) -- (258.0,329.0);
  \draw[line width=9.3pt] (223.0,311.0) -- (222.0,302.0);
  \draw[line width=3.5pt] (535.0,508.0) -- (528.0,508.0);
  \draw[line width=5.0pt] (829.9,191.1) -- (824.0,216.7);
  \draw[line width=4.0pt] (876.0,380.0) -- (910.0,368.0);
  \draw[line width=5.0pt] (1127.5,379.0) -- (1151.2,370.8);
  \draw[line width=5.0pt] (1084.8,119.8) -- (1033.0,95.0);

  %% ════ 曲线（prims2 的 curves —— 三段式，坐标同为 (y,x)）════
  \draw[line width=3.0pt] (611.0,509.0) .. controls (607.9,506.1) and (604.0,506.9) .. (601.0,510.0);
  \draw[line width=5.0pt] (225.0,248.0) .. controls (231.2,252.5) and (235.5,244.1) .. (230.0,240.0);
  \draw[line width=6.0pt] (985.0,225.0) .. controls (996.2,234.2) and (1010.2,245.7) .. (1024.9,237.0);
  \draw[line width=7.0pt] (1024.9,237.0) .. controls (1046.7,221.5) and (1059.0,189.0) .. (1088.0,188.0);
  \draw[line width=6.0pt] (983.0,224.0) .. controls (945.8,183.6) and (881.3,226.8) .. (876.0,272.0);
  \draw[line width=4.0pt] (73.0,217.0) .. controls (51.5,190.1) and (32.1,155.8) .. (40.0,121.1);
  \draw[line width=3.0pt] (40.0,121.1) .. controls (44.9,103.5) and (58.4,89.2) .. (73.7,79.3);
  \draw[line width=4.0pt] (73.7,79.3) .. controls (97.9,65.5) and (124.3,59.4) .. (150.8,53.0);
  \draw[line width=4.0pt] (150.8,53.0) .. controls (232.7,40.9) and (323.7,41.4) .. (405.1,52.0);
  \draw[line width=4.0pt] (415.7,51.0) .. controls (511.2,63.7) and (605.2,80.8) .. (695.8,109.0);
  \draw[line width=5.0pt] (442.0,423.0) .. controls (448.3,427.2) and (440.9,436.3) .. (436.0,430.0);
  \draw[line width=3.0pt] (529.0,501.0) .. controls (532.0,499.5) and (537.6,498.9) .. (540.0,502.0);
  \draw[line width=5.0pt] (1033.0,95.0) .. controls (955.9,69.8) and (844.9,102.5) .. (829.9,191.1);
  \draw[line width=4.0pt] (824.0,216.7) .. controls (803.2,265.2) and (789.8,336.6) .. (834.9,376.9);
  \draw[line width=5.0pt] (834.9,376.9) .. controls (847.3,383.5) and (862.5,382.4) .. (876.0,380.0);
  \draw[line width=5.0pt] (910.0,368.0) .. controls (980.3,324.7) and (1055.0,377.4) .. (1127.5,379.0);
  \draw[line width=5.0pt] (1151.2,370.8) .. controls (1163.8,361.8) and (1171.5,348.0) .. (1178.0,334.5);
  \draw[line width=5.0pt] (1178.0,334.5) .. controls (1201.4,253.7) and (1167.1,151.8) .. (1084.8,119.8);

  %% ════ 实心点 ════
  \fill (740.0,132.5) circle (3.4);
  \fill (345.8,171.8) circle (7.0);
  \fill (875.4,273.5) circle (7.0);
  \fill (440.5,422.5) circle (5.1);

  %% ════ 标签（probe 直接给的 \node 行）════
  %% ⚠ 全部补了 fill=white：文字在最上层，否则与曲线交叉干扰
     \node[anchor=center,inner sep=0pt,fill=white,font=\fontsize{28.5}{35.6}\selectfont] at (485.0,28.5) {\textsf{\textbf{\textit{f}}}};   % box(479,17,11,22)
     \node[anchor=center,inner sep=0pt,fill=white,font=\fontsize{31.4}{39.2}\selectfont] at (130.1,147.3) {\textsf{\textbf{\textit{u}}}};   % box(120,137,17,17)
     \node[anchor=center,inner sep=0pt,fill=white,font=\fontsize{40.2}{50.3}\selectfont] at (1102.7,222.5) {\textsf{\textbf{×}}};   % box(1092,211,20,17)
     \node[anchor=center,inner sep=0pt,fill=white,font=\fontsize{22.2}{27.8}\selectfont] at (1121.1,230.4) {\textsf{\textbf{1}}};   % box(1115,222,7,16)
     \node[anchor=center,inner sep=0pt,fill=white,font=\fontsize{24.3}{30.3}\selectfont] at (202.4,241.4) {\textsf{\textbf{\textit{f}}}};   % box(198,230,8,22)
     \node[anchor=center,inner sep=0pt,fill=white,font=\fontsize{30.7}{38.4}\selectfont] at (214.5,244.5) {\textsf{\textbf{(}}};   % box(210,230,8,28)
     \node[anchor=center,inner sep=0pt,fill=white,font=\fontsize{30.2}{37.7}\selectfont] at (245.5,244.0) {\textsf{\textbf{)}}};   % box(240,230,8,27)
     \node[anchor=center,inner sep=0pt,fill=white,font=\fontsize{30.5}{38.1}\selectfont] at (160.5,245.3) {\textsf{\textbf{\textit{u}}}};   % box(151,235,16,17)
     \node[anchor=center,inner sep=0pt,fill=white,font=\fontsize{30.3}{37.9}\selectfont] at (187.8,250.6) {{\mfallback\textbf{−}}};   % box(172,244,17,4)
     \node[anchor=center,inner sep=0pt,fill=white,font=\fontsize{29.2}{36.5}\selectfont] at (875.0,307.4) {\textsf{\textbf{\textit{x}}}};   % box(866,297,16,17)
     \node[anchor=center,inner sep=0pt,fill=white,font=\fontsize{24.1}{30.2}\selectfont] at (890.0,317.2) {\textsf{\textbf{0}}};   % box(882,309,12,16)
     \node[anchor=center,inner sep=0pt,fill=white,font=\fontsize{38.1}{47.6}\selectfont] at (987.8,417.5) {\textsf{\textbf{\textit{x}}}};   % box(976,404,21,22)
     \node[anchor=center,inner sep=0pt,fill=white,font=\fontsize{28.7}{35.9}\selectfont] at (214.0,475.5) {\textsf{\textbf{(}}};   % box(210,461,7,28)
     \node[anchor=center,inner sep=0pt,fill=white,font=\fontsize{27.2}{34.0}\selectfont] at (276.5,472.5) {\textsf{\textbf{(}}};   % box(272,461,8,22)
     \node[anchor=center,inner sep=0pt,fill=white,font=\fontsize{30.7}{38.4}\selectfont] at (304.5,475.5) {\textsf{\textbf{(}}};   % box(300,461,8,28)
     \node[anchor=center,inner sep=0pt,fill=white,font=\fontsize{30.2}{37.7}\selectfont] at (290.5,476.0) {\textsf{\textbf{)}}};   % box(285,462,8,27)
     \node[anchor=center,inner sep=0pt,fill=white,font=\fontsize{30.2}{37.7}\selectfont] at (336.5,476.0) {\textsf{\textbf{)}}};   % box(331,462,8,27)
     \node[anchor=center,inner sep=0pt,fill=white,font=\fontsize{31.4}{39.2}\selectfont] at (174.1,476.3) {\textsf{\textbf{\textit{u}}}};   % box(164,466,17,17)
     \node[anchor=center,inner sep=0pt,fill=white,font=\fontsize{31.4}{39.3}\selectfont] at (231.8,476.1) {\textsf{\textbf{\textit{e}}}};   % box(222,466,16,17)
     \node[anchor=center,inner sep=0pt,fill=white,font=\fontsize{29.0}{36.3}\selectfont] at (322.0,476.5) {\textsf{\textbf{\textit{a}}}};   % box(314,467,14,16)
     \node[anchor=center,inner sep=0pt,fill=white,font=\fontsize{24.9}{31.1}\selectfont] at (246.1,485.7) {\textsf{\textbf{0}}};   % box(238,477,12,17)
     \node[anchor=center,inner sep=0pt,fill=white,font=\fontsize{27.5}{34.3}\selectfont] at (675.2,506.3) {\textsf{\textbf{\textit{z}}}};   % box(667,498,15,13)
     \node[anchor=center,inner sep=0pt,fill=white,font=\fontsize{23.8}{29.8}\selectfont] at (632.1,506.2) {\textsf{\textbf{\textit{c}}}};   % box(625,499,13,12)

  % 钉死画布：否则超出的图元/控制点会把页面撑大，1:1 回比就失准
  \pgfresetboundingbox
  \path[use as bounding box] (0,0) rectangle (1207,511);
\end{tikzpicture}
\end{document}

% ⚠ 待核清单（probe 拿不准，人眼过一遍）：
%   · 未识别/跳过：⑤ 跳过/未识别的字形
%   · 未识别/跳过：box(254,465,12,13)
